\section{Proof Details}
\label{app:proofs}

\subsection{Boundary-complete single-crossing proof}

For \eqref{eq:exp-linear}, direct differentiation of
\eqref{eq:objective} gives
\[
 J'(b)=\frac{A(b)}{[1-qp(b)]^2},
\]
where \(A(b)\), \(K\), and \(g(b)\) are defined in
\eqref{eq:Ab}.  The denominator is positive.  The slope of \(g\) is
\(q\bar p\alpha\beta C_t>0\).

If \(g(0)\geq0\), then \(g(b)>0\) for every \(b>0\), so
\(A(b)>K>0\).  If \(g(0)<0\), let \(b_g>0\) be its unique zero.  On
\([0,b_g)\), putting \(h(b)=e^{-\alpha b}g(b)\) yields
\[
 h'(b)=e^{-\alpha b}[g'(b)-\alpha g(b)]>0.
\]
Thus \(A=K+h\) increases strictly until \(b_g\), while
\(A(b)\geq K>0\) for \(b\geq b_g\).  When \(A(0)>0\), the derivative is
positive everywhere.  When \(A(0)=0\), it is zero only at the left endpoint
and positive for all \(b>0\).  When \(A(0)<0\), continuity and strict
increase on \([0,b_g]\) give one negative-to-positive crossing.  This proves
all shape and boundary claims in Theorem~\ref{thm:single-crossing}.

Finally,
\begin{align*}
 A(0)<0
 &\Longleftrightarrow
 q\bar p\alpha[(1-q)C_v-HC_t]>\beta C_t\\
 &\Longleftrightarrow
 R>\frac{H+\beta/(q\bar p\alpha)}{1-q}
 =\frac{H+\delta/q}{1-q}.
\end{align*}

\subsection{Lambert-W derivation}

In the active region,
\[
 A(b)=K+q\bar p(c_0+c_1b)e^{-\alpha b}=0.
\]
Let \(x=b+c_0/c_1\) and \(w=-\alpha x\).  Then
\[
 x e^{-\alpha x}
 =-\frac{K}{q\bar p c_1}e^{-\alpha c_0/c_1}
\]
and hence
\[
 w e^w=
 \frac{\alpha K}{q\bar p c_1}
 e^{-\alpha c_0/c_1}=\theta>0.
\]
The real solution is \(w=W_0(\theta)>0\), yielding
\eqref{eq:lambert}.  Theorem~\ref{thm:single-crossing} guarantees that this
root is the unique unbounded interior minimizer.  Restricting the domain to
\([0,b_{\max}]\) gives \eqref{eq:clipping} by the decreasing-then-increasing
shape.

\subsection{Interior comparative statics}

At an interior optimum, \(A_b>0\) by the single-crossing proof, so for a
parameter \(x\),
\[
 \frac{d b_\infty^\star}{dx}=-\frac{A_x}{A_b}.
\]
Holding the other primitive parameters fixed,
\begin{align*}
 A_{C_v}&=q p'(q-1)<0,\\
 A_{C_t}&=q p'(G+\tau)+\tau'(1-qp)>0,\\
 A_G&=q p'C_t>0.
\end{align*}
For the exponential--linear family,
\[
 A_\beta=C_t[q p'b+(1-qp)]>0.
\]
Therefore \(b_\infty^\star\) increases with \(C_v\) and decreases with
\(C_t\), \(G\), and \(\beta\), until clipping turns strict changes into weak
ones.

\subsection{Participation interval and strict unimodality}

Differentiation gives
\[
 F'(q)=
 \frac{Hq^2+\delta(2q-1)}{q^2(1-q)^2}.
\]
The numerator has derivative \(2Hq+2\delta>0\), starts negative, and ends
positive.  Solving its unique zero and evaluating \(F\) there gives the
\(q_0\) and \(F_{\min}\) in Theorem~\ref{thm:q-shape}.

The equality \(R=F(q)\) is equivalent to
\[
 Rq^2-(R-H)q+\delta=0.
\]
When \(R>F_{\min}\), its two roots are \eqref{eq:qroots} and lie in
\((0,1)\); strict U-shape then makes \(R>F(q)\) hold exactly between them.
When \(R\leq F_{\min}\), the strict participation inequality has no solution
in \((0,1)\).  The quadratic discriminant can also be nonnegative on a
low-\(R\) algebraic branch, but those roots are nonpositive and do not create
a feasible participation interval.

For the shape inside \((q_-,q_+)\), differentiation of \(A\) with respect to
\(q\), followed by collection of terms, gives
\begin{align}
 A_q(b,q)&=2C_v p'(b)[q-q_c(b)],\\
 q_c(b)&=\frac12\left[
 1-\frac{H+\beta b-\frac{\beta}{\alpha}(e^{\alpha b}-1)}{R}
 \right].
\end{align}
At the interior optimum \(A_b>0\).  The implicit-function theorem gives
\[
 \operatorname{sgn}\frac{d b_\infty^\star}{dq}
 =-\operatorname{sgn}\Psi(q),\qquad
 \Psi(q):=q-q_c(b_\infty^\star(q)).
\]
At any zero \(\widehat q\) of \(\Psi\), the preceding derivative vanishes,
and therefore
\[
 \Psi'(\widehat q)
 =1-q_c'(b_\infty^\star)
 \frac{db_\infty^\star}{dq}\bigg|_{\widehat q}=1.
\]
Every zero is a strict upward crossing.  Two such zeros would require a
downward crossing between them, so there is at most one.  The optimum is
positive inside the active interval and tends to zero at both endpoints;
therefore it has an interior maximum and at least one zero of \(\Psi\).
The zero is unique, with \(\Psi<0\) before it and \(\Psi>0\) after it.  This
proves strict increase followed by strict decrease.

\subsection{Strong-order proof for the static index}

Let
\(\varphi(b,Q)=C_vV(b)+Q S_{\rm eff}(b)\).  Suppose
\(Q_2>Q_1\), \(b_i\in\mathcal M(Q_i)\), and, toward a contradiction,
\(b_2>b_1\).  Optimality gives
\begin{align*}
 \varphi(b_1,Q_1)&\leq\varphi(b_2,Q_1),\\
 \varphi(b_2,Q_2)&\leq\varphi(b_1,Q_2).
\end{align*}
Adding and canceling the risk terms yields
\[
 (Q_2-Q_1)
 [S_{\rm eff}(b_2)-S_{\rm eff}(b_1)]\leq0,
\]
which contradicts strict increase of \(S_{\rm eff}\).  Hence every high-\(Q\)
minimizer is no larger than every low-\(Q\) minimizer.

\subsection{Why history-invariant detectability matters}
\label{app:persistence}

Marginal values of \(q\) and \(p\) do not suffice for
Proposition~\ref{prop:renewal}.  For example, set \(q=p=1/2\), but draw once
per intent a latent detector type: with probability \(1/2\) it catches every
unsafe draft, and with probability \(1/2\) it catches none.  Draft safety is
otherwise independent across rounds.  Conditional on the perfect detector,
exit waits for a safe draft, so \(\E[N]=2\) and \(V=0\); conditional on the
blind detector, exit occurs in one round and \(V=1/2\).  The mixture has
\[
 \E[N]=\frac32,\qquad V=\frac14,
\]
whereas substituting the same marginals into the independent-pair formulas
would give \(4/3\) and \(1/3\).  Persistent blind spots must therefore be
modeled through a latent-state or history-dependent process.
